\scanpage{46}
% loc: heading, lines 1–2
\textcircled{$\alpha$}. \uline{Foncteurs pleinement fidèles sur catégories des motifs}, \uline{liés à théorie de Hodge.}

% loc: opening, lines 3–5
$S$ un $\mathbb{Q}$-schéma régulier\inserted{ connexe}, muni d'un point géom $\unsure{\xi|\zeta}:\operatorname{Spec}(\mathbb{C})\longrightarrow S$
Le foncteur $\xi^*:\underset{\text{\inserted{motifs lisses sur $S$}}}{\operatorname{Mot}(S)}\longrightarrow\operatorname{Mot}(\mathbb{C})$ est fidèle. On \unsure{regardera|regarde} le foncteur composé
% The hook below Mot(S) attaches the gloss to the source.
% loc: diagram below the opening
\diagram{\begin{tikzcd}[ampersand replacement=\&,column sep=2.1em]
\operatorname{Mot}(S) \arrow[r,"\xi^*"] \arrow[rrr,bend right=14,"{(T_{\mathrm B})_\xi}"'] \& \operatorname{Mot}(\mathbb C) \arrow[r,"{T_{\deleted{\unsure{\mathrm B}}\mathrm{B+H}}}"] \& \operatorname{Hodge}(\mathbb Q,\mathbb C) \arrow[r] \& \operatorname{Modf}(\mathbb Q)
\end{tikzcd}}

% loc: left of item 1
% slanted
\margin{\begin{tikzpicture}[baseline=(n.base)]\node[draw,ellipse,inner sep=2pt] (n) {réseau et vectoriel filtré};\end{tikzpicture}}
% loc: item 1, first line
\textcircled{1} \inserted{(Via Hodge)} Supposons $\xi$ localisé au $\mathrm{p}^{\mathrm t}$ générique $\unsure{y|\eta}$ de $S$. \deleted{Le fonct} \textsuperscript{\textcircled{a}} \unsure{Á|À} tout motif $M$ sur $S$, on associe le triple
% loc: first triple
\[
\bigl\{T_{\mathrm{DR}}(M_y),T_{\mathrm B}(M_\xi),\varphi_M\bigr\}
\]
% loc: first comparison, line below triple
où $\varphi_M$ est l'isom can $\varphi_M:T_{\mathrm B}(M_\xi)\otimes_{\mathbb Q}\mathbb C\simeq T_{\mathrm{DR}}(M_y)\otimes_{k(y)}\mathbb C$.

% loc: slanted question by the two comparisons
% slanted
\margin{Le foncteur reste-t-il pl. fidèle si on laisse tomber la filtration de $T_{\mathrm{DR}}$ ?? id pour bigr. de $T_{\mathrm{Hdg}}$ !!}
% loc: filtration paragraph
On considère $T_{\mathrm{DR}}(M_y)$ comme \uline{filtré}, la filtration étant telle que $\varphi_M$ la transforme en une filtration \uline{admissible} du complexifié de $T_{\mathrm B}(M_\xi)$ [i.e. définissant, avec la filtration complexe conjuguée, une bigraduation].

% loc: second functor, lines above second triple
Le foncteur \inserted{en $M$} ainsi obtenu est pleinement fidèle (moyennant conj. de Hodge simple).\textsuperscript{\textcircled{b}} Même résultat si on considère
% loc: second triple
\[
\bigl\{T_{\mathrm{Hdg}}(M),T_{\mathrm B}(M_\xi),\Psi_M\bigr\}
\]
% loc: left of second triple
% slanted
\margin{\begin{tikzpicture}[baseline=(n.base)]\node[draw,ellipse,align=left,inner sep=3pt] (n) {réseau et\\vectoriel\\bigradué};\end{tikzpicture}}
% loc: final comparison
$T_{\mathrm{Hdg}}(M)$ étant considéré comme \uline{bigradué}, et $\Psi_M$ étant l'isomorphisme canonique $\Psi_M:T_{\deleted{\unsure{\mathrm{DR}|\mathrm B}}\mathrm B}(M_\xi)\unsure{\otimes|\ominus}_{\deleted{\unsure{k|\mathbb Q}}\mathbb Q}\mathbb C\simeq T_{\unsure{\mathrm{Hgg}|\mathrm{Hdg}}}(M)\otimes_{k(y)}\mathbb C$.
% loc: paper watermark, read sideways and compared across sheets
% printed form
% paper watermark, visible letters in the direction of the paper lettering
\par\noindent\textsc{trong}\hfill\textsc{Extra}\\
\unsure{ry|y} 1808\hfill\unsure{Maunu|Mauno|Maun}
% archive sheet number: 45

